\chapter{Introduction}

\section{Contexte du module}
Dans ce module, je travaille sur le traitement de nuages de points 3D.
Un point cloud est un ensemble de points \((x,y,z)\) non ordonnes: je ne peux donc pas supposer un voisinage fixe comme dans une image 2D.
Le probleme principal est de construire un modele qui respecte cette structure en ensemble.

\section{Lien avec l'article PointNet}
Le projet est inspire de l'article \textit{PointNet: Deep Learning on Point Sets for 3D Classification and Segmentation}~\cite{qi2017pointnet}.
Avant d'implementer mes reseaux, j'ai essaye de comprendre sa logique:
\begin{itemize}
  \item traiter directement les points, sans voxelisation lourde;
  \item apprendre des descripteurs locaux par point via des MLP partages;
  \item utiliser une fonction symetrique (max-pooling) pour obtenir une invariance a l'ordre des points;
  \item ajouter un T-Net pour stabiliser l'alignement geometrique en entree.
\end{itemize}

\section{Objectif de mon rapport}
Je veux montrer que je comprends ce que je fais, et pas seulement que j'obtiens un score.
Concretement, je documente:
\begin{itemize}
  \item mon evolution de PointMLP vers PointNet;
  \item ma methode experimentale (runs, controles, comparaison);
  \item pourquoi j'ai retenu \texttt{baseline + occlusion} en fin de projet.
\end{itemize}
